% Generated draft: requires mathematical and source review.
\documentclass[11pt]{article}
\usepackage[margin=25mm]{geometry}
\usepackage{fontspec}
\setmainfont{Latin Modern Roman}
\usepackage{amsmath,amssymb,mathtools,unicode-math}
\setmathfont{Latin Modern Math}
\usepackage{enumitem,needspace,xurl,hyperref}
\setlength{\parindent}{0pt}
\setlength{\parskip}{5pt}
\setlength{\emergencystretch}{3em}
\newcommand{\N}{\mathbb{N}}
\newcommand{\Z}{\mathbb{Z}}
\newcommand{\Q}{\mathbb{Q}}
\newcommand{\R}{\mathbb{R}}
\newcommand{\C}{\mathbb{C}}
\newcommand{\F}{\mathbb{F}}
\newcommand{\A}{\mathbb{A}}
\newcommand{\PP}{\mathbb{P}}
\newcommand{\E}{\mathbb{E}}
\newcommand{\eps}{\varepsilon}
\newcommand{\dd}{\,\mathrm{d}}
\begin{document}
\section*{UnsolvedMath ID: 20000833}
\textbf{Source problem code:} catalogue entry
\subsection{Definitions and mathematical statement}
Let \(H\) be an irreducible component of \(\operatorname{Hilb}_{d,g}(\PP_k^3)\), with Hilbert polynomial \(dt+1-g\). Interpret general points over an algebraically closed characteristic-zero field \(k\) as tuples in a dense open subset of \((\PP^3)^r\), rather than merely a finite list of linear-position conditions. Define
\[
 I_r(H)=\{([C],p_1,\ldots,p_r):[C]\in H,\ p_i\in C,
 \ p_i\ne p_j\ (i\ne j)\},
\]
\[
 \iota(H)=\max\{r\ge0:I_r(H)\longrightarrow(\PP^3)^r
 \text{ is dominant}\}.
\]
The source asks to determine \(\iota(H)\) for fixed \(d,g\) and a specified component \(H\). It does not specify a field or the exact meaning of general position; the above are explicit interpretation choices. Its remark says the problem is hard but should be treatable for arithmetically Cohen--Macaulay curves, meaning curves whose homogeneous coordinate rings are Cohen--Macaulay.

For a smooth curve \(C\) at a smooth Hilbert point and distinct marked smooth points, the supplied aid gives the normal-evaluation criterion
\[
 H^0(C,N_{C/\PP^3})\longrightarrow
 \bigoplus_{i=1}^r N_{C/\PP^3}|_{p_i}.
\]
Surjectivity is the infinitesimal criterion for the incidence projection to be a submersion at that flag. In characteristic zero, dominance on the relevant incidence component is tested at its general smooth flags; dominance is not a claim that evaluation is surjective at every special flag. This criterion applies to components with the necessary smooth-curve locus and is not a complete answer for arbitrary Hilbert components.

The aid reports an exact component-wide answer for the complete-intersection component \(H_{a,b}\), the closure of smooth complete intersections of surfaces of degrees \(1\le a\le b\). Here
\[
 d=ab,\qquad g=1+\frac{ab(a+b-4)}2,
\]
and the reported answer is
\[
 \iota(H_{a,b})=
 \begin{cases}
 \binom{a+3}{3}-1,&a<b,\\
 \binom{a+3}{3}-2,&a=b.
 \end{cases}
\]
The supplied lower-bound argument uses
\(N_{C/\PP^3}\simeq\mathcal O_C(a)\oplus\mathcal O_C(b)\)
and normal evaluation. For the upper bound, every member of this component, including non-complete-intersection boundary subschemes, is reported to retain at least one independent degree-\(a\) equation if \(a<b\), and at least two if \(a=b\), by upper semicontinuity. Thus the formula is asserted for the whole component, not only its smooth complete-intersection open locus. The record does not provide a corresponding formula for every arithmetically Cohen--Macaulay component or every pair \((d,g)\).
\subsection{Short English statement}
For a chosen component of space curves, how many general points can some member contain? The supplied aid gives a normal-bundle evaluation test and a sharp formula for complete-intersection components, including a claimed upper bound that also controls their boundary curves.
\subsection{Sources}
\begin{itemize}
\item Local UnsolvedMath registry, permanent ID \(20000833\), source code \texttt{catalogue entry}; \texttt{aim-arithmetic-geometry-notes.json}, \texttt{notes[78]}, tag \texttt{problem}.
\item AIM workshop ``Components of Hilbert schemes,'' Problem 32: \url{https://aimath.org/WWN/hilbertschemes/hilbertschemes.pdf}. Statement status \texttt{exact}; original remark calls the problem hard and suggests the arithmetically Cohen--Macaulay case should be treatable. No original literature citation, proposer, or year is supplied. No PDF was accessed for this draft.
\item Research attempt 1: \texttt{partial\_result}, problem status at run \texttt{open}, rigor label \texttt{proved}, novelty confidence low; registry status \texttt{partially\_solved}, classification \texttt{PARTIAL-PROGRESS}. Registry created and updated 14 August 2026. No authoritative external citation is supplied for the proposed exact formula. This machine-generated aid is not peer reviewed and requires independent verification before citation.
\item Ulam AI / UnsolvedMath Contributors; CC BY 4.0.
\end{itemize}
\end{document}
